\documentclass[11pt]{article}
\usepackage[margin=0.9in]{geometry}
\usepackage{amsmath,array,booktabs,enumitem,tikz}
\usepackage[version=4]{mhchem}
\usetikzlibrary{arrows.meta}
\setlength{\parindent}{0pt}\setlength{\parskip}{6pt}
\renewcommand{\arraystretch}{1.5}
\pagestyle{plain}
% orbital box: 0 empty, 1 up, 2 paired
\newcommand{\ob}[1]{\tikz[baseline=-0.7ex]{\draw (0,-0.28) rectangle (0.5,0.28);%
\ifcase#1\or\draw[-{Stealth[length=3pt]}](0.25,-0.2)--(0.25,0.2);%
\or\draw[-{Stealth[length=3pt]}](0.17,-0.2)--(0.17,0.2);\draw[-{Stealth[length=3pt]}](0.33,0.2)--(0.33,-0.2);\fi}}
\newcommand{\sub}[2]{\begin{tabular}[b]{@{}c@{}}#2\\[-6pt]\small #1\end{tabular}\hspace{10pt}}
\newcommand{\pb}[3]{\ob{#1}\hspace{-0.4pt}\ob{#2}\hspace{-0.4pt}\ob{#3}}
\newcommand{\ans}[1][1]{\par\foreach \i in {1,...,#1}{\vspace{0.8cm}\hrule}\vspace{4pt}}
\newcommand{\nm}{Name: \rule{7cm}{0.4pt}\par}

\begin{document}
{\Large\bfseries Electrons, Ions and Groups}\par\nm
\section*{1. Electron configuration}
Electrons fill the lowest energy levels first, in this order:
\[ 1s \rightarrow 2s \rightarrow 2p \rightarrow 3s \rightarrow 3p \rightarrow 4s \]
An $s$ sublevel holds 2 electrons. A $p$ sublevel holds 6.

\textbf{Example: sulfur (16 electrons).} \quad $1s^2\,2s^2\,2p^6\,3s^2\,3p^4$ \quad $(2+2+6+2+4=16)$
\begin{itemize}[nosep]
\item \textbf{Energy levels used:} the highest number in the configuration. For sulfur, 3.
\item \textbf{Valence electrons:} the electrons in the outer level. For sulfur, $3s^2\,3p^4 = 6$.
\item \textbf{Core electrons:} all the rest. For sulfur, $16-6=10$.
\end{itemize}

\textbf{Orbital boxes.} Each box holds 2 electrons with opposite arrows. In a $p$ sublevel, put one electron in each of the three boxes before pairing any.

Nitrogen (7 electrons):\quad \sub{1s}{\ob2}\sub{2s}{\ob2}\sub{2p}{\pb111}

\textbf{Try it.} Draw the orbital boxes for carbon (6 electrons).\par\vspace{6pt}
\sub{1s}{\ob0}\sub{2s}{\ob0}\sub{2p}{\pb000}

\section*{2. Cations and anions}
An ion is an atom that has gained or lost electrons. The number of protons never changes.

\begin{tabular}{llll}\toprule
 & Electrons & Charge & Name\\\midrule
Lose electrons & fewer than protons & positive ($+$) & cation\\
Gain electrons & more than protons & negative ($-$) & anion\\\bottomrule
\end{tabular}

\textbf{Example: \ce{Mg^2+}.} Magnesium has 12 protons. A $2+$ charge means 2 fewer electrons than protons, so 10 electrons. Cation.

\textbf{Example: \ce{Cl-}.} Chlorine has 17 protons. A $1-$ charge means 1 more electron than protons, so 18 electrons. Anion.

\textbf{Try it.} How many electrons does \ce{O^2-} have? Is it a cation or an anion?\ans

\section*{3. Groups of the periodic table}
Elements in the same group have the same number of valence electrons, so they behave the same way. Atoms gain or lose electrons to get a full outer level.

\begin{tabular}{clccp{5.2cm}}\toprule
Group & Name & Valence $e^-$ & Ion & What they do\\\midrule
1 & Alkali metals & 1 & $1+$ & Lose 1 electron. The most reactive metals.\\
2 & Alkaline earth metals & 2 & $2+$ & Lose 2 electrons. Reactive, but less than group 1.\\
17 & Halogens & 7 & $1-$ & Gain 1 electron. The most reactive nonmetals.\\
18 & Noble gases & 8 & none & Outer level already full. Almost never react.\\\bottomrule
\end{tabular}

Hydrogen sits in group 1 but is not an alkali metal; it is a nonmetal gas.

\textbf{Example: which is more reactive, sodium or magnesium?} Sodium is in group 1 and loses one electron. Magnesium is in group 2 and must lose two. Losing one takes less energy, so sodium is more reactive.

\textbf{Try it.} Which is more reactive, potassium or calcium? Why?\ans[2]
\end{document}
